\documentclass[10pt]{article}

\usepackage[utf8]{inputenc}

\usepackage[T1]{fontenc}

\usepackage{amsmath}

\usepackage{amsfonts}

\usepackage{amssymb}

\usepackage[version=4]{mhchem}

\usepackage{stmaryrd}

\usepackage{bbold}



\begin{document}

Голоморфная функция $f$ бесконечно дифференцируема, ряд (10) есть ее ряд Тейлора, коэффициенты к-рого выражаются формулами



\[

c_{k}=\left.\frac{1}{k_{1} ! \ldots k_{n} !} \frac{\partial^{k_{1}+\ldots k_{n}}}{\partial z_{1}^{k_{1}} \ldots \partial z_{n}^{k_{n}}}\right|_{a}

\]



через значения производных функции $f$ в точке $a$.



Основная интегральная теорема Коши-Пуа н каре для голоморфной функщии $f$ в области $D \subset \mathbb{C}^{n}$ имеет вид:



\[

\int_{\partial G} f d z=0

\]



где $G-$ произвольная $(n+1)$-мерная поверхность, $G \subset \bar{G} \subset D, \quad \delta G-$ кусочно гладкая граница $G$; в плоском случае $n+1=2 n$, а при $n>1 \quad n+1<2 n$. Аналог интегральной формулы Коши просто выписывается для поликруговой области $D=D_{1} \times \ldots \times D_{n}, D_{k}$ - плоская область:



\[

f(z)=\frac{1}{(2 \pi i)^{n}} \int_{\Gamma} \frac{f(\zeta) d \zeta}{\zeta-z} .

\]



Здесь $\Gamma=\partial D_{1} \times \ldots \times \partial D_{n}-n$-мерный остов границы $\partial D$, $\zeta=\left(\zeta_{1}, \ldots, \zeta_{n}\right)$



\[

\frac{d \zeta}{\zeta-z}=\frac{d \zeta_{1} \ldots d \zeta_{n}}{\left(\zeta_{1}-z_{1}\right) \ldots\left(\zeta_{n}-z_{n}\right)}

\]



Однако в областях общего вида разделение переменных невозможно и интегральные формулы усложняются.



Изучение особенностей А. ф. нескольких переменных представляет значительный интерес. Основное отличие от плоского случая состоит в следующем (т е о е м а О с гу да-Б р а у на): если $f-$ голоморфная функция в $D \backslash K$, где $D$ - область в $\mathbb{C}^{n}, n>1, K-$ компакт в $D$, не разбивающий $D$, то $f$ голоморфно продолжается во всю область $D$. Следовательно, голоморфные функции $f$ нескольких переменных не могут иметь изолированных особенностей. Поэтому основная задача многомерной теории вычетов значительно сложнее одномерной. Она состоит в вычислении интеграла



\[

\int_{\sigma} f d z

\]



от функции $f$, голоморфной в области $D$, кроме (аналитического) множества $M$, по значениям $f$ на $n$-мерной поверхности $\sigma$, не пересекающейся с $M$. Если $\sigma$ не зацеплена с $M$, то есть ограничивает $(n+1)$-мерную поверхность $G$ такую, что $G \subset \bar{G} \subset D \backslash M$, то по теореме Коши-Пуанкаре



\[

\int_{\sigma} f d z=0

\]



В общем случае нужно выяснить характер зацепления $\sigma$ с $M$ и вычислить интегралы по специальным $n$-мерным поверхностям, определяемым множеством $M$, к-рые теперь играют роль вычетов (см. [3], [4]).



Какова бы ни была область $D \subset \mathbb{C}$, можно построить А. ф. $f$, к-рая голоморфна в $D$ и не продолжается аналитически вне $D$, то есть $D$ является для $f$ естественной областью сушествования. В случае $D \subset \mathbb{C}^{n}, n>1$, может оказаться, что $D$ не является областью голоморфности, то есть любая А. ф. в области $D$ продолжается голоморфно за пределы $D$.



Для того чтобы область $D \subset \mathbb{C}^{n}$ была областью голоморфности, необходимо и достаточно, чтобы она была псевдовыпуклой. Последнее означает, что действительная функция $u(z)=-\ln d(z ; \partial D)$ должна быть плюрисубгармонической (см. [4]). Здесь $d$ - евклидово расстояние точки $z \in D$ до границы $\partial D$. Наконец, для того чтобы действительная функция $u(z)$ класса $\mathbb{C}^{2}(D), u(z)<+\infty$, была плюрисубгар- монической, необходимо и достаточно, чтобы форма Л е в и



\[

\boldsymbol{H}_{z}(u ; \omega)=\left.\sum_{\mu, v=1}^{n} \frac{\partial^{2} u}{\partial z_{\mu} \partial \bar{z}_{v}}\right|_{z} \omega_{\mu} \omega_{v}

\]



была неотрицательно определенной при всех $\omega \in \mathbb{C}^{n}, z \in D$. Простейшее достаточное условие области голоморфности $D$ состоит в том, что $D$ должна быть выпуклой.



Основной в этом круге вопросов является следующая задача: если область $D$ не является областью голоморфности, то требуется построить ее оболочку голоморфности, то есть наименьшую область голоморфности, в к-рую аналитически продолжается любая функция, голоморфная в $D$. При этом следует иметь в виду, что при аналитич. продолжении (как и в плоском случае) область голоморфности $\begin{array}{ll}\text { А. ф. может оказаться многозначной }\end{array}$ (см. [4]).



Важное значение (как и в плоском случае) имеют задачи построения в данной области голоморфной с заданным нулями или мероморфной функции с заданными главными частями особенностей (п робл е ма Кузена). В плоском случае эти задачи решаются в произвольной области, согласно теоремам Вейерштрасса и Миттаг-Леффлера. В областях $D \subset \mathbb{C}^{n}, n>1$, требуется выполнение нек-рых дополнительных условий (см. [4]).



Лит.: [1] П р и в ало в И. И., Введение в теорию функций комплексного переменного, 2 изд., М., 1967; [2] Мар куш е в ич А. И., Теория аналитических функший, 2 изд., М., 1967; [3] В лад и м и ро в В. С., Методы теории функций многих комплексных переменных, М., 1964; [4] Ша бат Б. В., Введение в комплексный анализ, $\begin{array}{ll}\text { ч. 1-2, } 3 \text { изд., М., } 1985 . & \text { Е. Д. Соломенцев. }\end{array}$



АНАЛИтИЧЕСКИЙ ВЕКТОР в пространстве представлен ий $T$ группы Ли $G-$ вектор $\xi \in V$, для котоporo отображение $g \rightarrow T(g) \xi$ является вещественно аналитической вектор-функцией на $G$ со значениями в $V$. Если $T$ - слабо непрерывное представление группы Ли $G$ в банаховом пространстве $V$, то множество $V^{\omega}$ аналитич. векторов плотно в $V$ (см. [1] - [3]). Существует обобщение этой теоремы на широкий класс представлений в локально выпуклых пространствах (см. [5]). Было доказано (см. [6]), что представление связной группы Ли $G$ в банаховом пространстве $V$ однозначно восстанавливается по соответствующему представлению алгебры Ли группы $G$ в пространстве $V^{\omega}$.



Аналитический вектор для неограничен н ого опе рат ора $A$ в банаховом пространстве $V$ с областью определения $D(A)$ определяется как вектор



\[

\xi \in \bigcap_{n=1}^{\infty} D\left(A^{n}\right)

\]



для к-рого ряд



\[

\sum \frac{x^{n}}{n !}\left\|A^{n} \xi\right\|

\]



имеет положительный радиус сходимости. Это понятие, введенное в [2], является частным случаем общего понятия А. в. - здесь в качестве группы Ли $G$ выступает множество точек действительной прямой с операцией сложения. Оно оказалось полезным в теории операторов в банаховых пространствах и в теории эллиптич. дифференциальных операторов.



Jum.: [1] Cartier P., Dixmier J., «Amer. J. Math.», 1958, v. 80, p. 131-45; [2] Нелс о н Э., «Математика», 1962, т. 6, № 3, с. 89131; [3] Гординг Л., «Математика», 1965, т. 9, № 5, с. 78-94; [4] Carti e r P., Vecteurs analytiques. Séminaire Bourbaki, Mai 1959; [5] Moore R. T., «Mem. Amer. Math. Soc.», 1968, v. 78; [6] Harish-Chandra, "Trans. Amer. Math. Soc.», 1953, v. 75, p. 185-243. А. А. Кириллов.



АНАЛИТИЧЕСКИЙ ВОЛНОВОЙ ФРОНТ - см. ВолнОВой фронт.



АНАЛИТИЧЕСКИЙ ФУНКЦИОНАЛ (в Широком с м ысле) - линейный непрерывный функционал на основном пространстве (см. Обобщенная функция), состоящем из аналитических функций. А. ф. различных классов используются в нелокальной аксиоматич. квантовой теории поля, в теории псевдодифференциальных операторов (диф-





\end{document}